% META: {"id":"1SPE-SECDEG-CO-048","chapitre":"1SPE-SECOND-DEGRE","type_objet":"corrige","exercice_ref":"1SPE-SECDEG-EX-048","capacites_codes":["C6"],"sources_inspiration":[],"status":"approved","origin":"generated","status_history":["generated","audited","accepted","approved"],"release_acceptance":"RELEASE_OWNER_BATCH_ACCEPTANCE","acceptance_closure_digest":"sha256:9b3ccf9a81c5520fbb7e03b7b2d3e7bf2057a02834b6d908bf3be5f49f3b553f","acceptance_receipt_digest":"sha256:4fad86f0282e0fa4e0419cca1ad6379ae7af5a280c04a4a90880f90a1c044242"}
% BEGIN-VERIFY
% from sympy import symbols, Rational, expand
% x = symbols('x')
% A = x*(20 - x)
% assert expand(A) == -x**2 + 20*x
% xs = Rational(-20, 2*(-1))
% assert xs == 10
% assert A.subs(x, xs) == 100
% END-VERIFY
\begin{corrige}{1SPE-SECDEG-EX-048}

\textbf{1. Expression de l'aire.}

\commentaireMarge{Le périmètre du rectangle vaut $2x + 2\ell = 40$.}

Si $x$ est la largeur, le périmètre s'écrit $2x + 2\ell = 40$, d'où $\ell = 20 - x$ (la longueur).

L'aire de l'enclos est donc :
\[
A(x) = x \times \ell = x(20 - x).
\]

\medskip
\textbf{2. Forme développée et sommet.}

En développant :
\[
A(x) = x(20 - x) = -x^2 + 20x.
\]
C'est bien un polynôme du second degré avec $a = -1$, $b = 20$, $c = 0$.

\commentaireMarge{$\alpha = -\dfrac{b}{2a}$.}

\[
\alpha = -\frac{b}{2a} = -\frac{20}{2 \times (-1)} = 10, \qquad \beta = A(10) = 10 \times (20 - 10) = 100.
\]
Le sommet est $S(10\,;\,100)$.

\medskip
\textbf{3. Dimensions optimales.}

Comme $a = -1 < 0$, le sommet correspond à un maximum. L'aire est donc maximale pour $x = 10$~m (largeur), soit une longueur $\ell = 20 - 10 = 10$~m : l'enclos optimal est un \textbf{carré} de côté $10$~m, d'aire maximale $100$~m$^2$.

\end{corrige}
